\begin{table}[!htbp]
\centering
\bicaption{关键比较的全程 DSR 配对差与置信区间}{Paired full-horizon DSR differences and confidence intervals}
\label{tab:ch5-dynamic-effects}
\begin{tabular}{lcrrr}
\toprule
比较对象 & 检验族 & 差值 & 单比较 $95\%$ 区间 & $N_{\mathrm{comp}}=31$ 区间 \\
\midrule
联合微调相对联合蒸馏 & 主 & $+0.70$ & $[+0.51, +0.88]$ & $[+0.39, +0.98]$ \\
联合微调相对去耦微调 & 主 & $+0.33$ & $[-0.08, +0.84]$ & $[-0.29, +1.20]$ \\
联合微调相对无期望微调 & 主 & $+0.11$ & $[-0.15, +0.41]$ & $[-0.28, +0.66]$ \\
微调 MPC 相对在线规划 & 主 & $+0.75$ & $[+0.41, +1.16]$ & $[+0.25, +1.42]$ \\
联合微调相对在线规划 & 补 & $+0.18$ & $[-0.17, +0.61]$ & $[-0.31, +0.91]$ \\
微调 MPC 相对蒸馏 MPC & 补 & $+0.21$ & $[-0.06, +0.48]$ & $[-0.23, +0.60]$ \\
\bottomrule
\end{tabular}
\sysutabnote{单位均为百分点（pp）。“主”表示预设八项主比较族之一，“补”表示补充机制比较；最后一列采用全部 $N_{\mathrm{comp}}=31$ 项 DSR 比较的 Bonferroni 同时置信区间。置信区间由 $10{,}000$ 次以工作负载为聚类单位的分块 Bootstrap 重采样计算，条件于三组固定上游预训练模型。}
\end{table}
